% Source: https://artofproblemsolving.com/community/c3h1530737_quadratics
% Stable identity: aops-topic-1530737-q01
Let her speed over the first $5$ km be $v$ km per hour, so the last $4$ km took place at $v-2>0$. The total time is one hour:
\[
\dfrac{5}{v}+\dfrac{4}{v-2}=1 .
\]
Multiplying by $v(v-2)$: $5(v-2)+4v=v^2-2v$, i.e.\ $v^2-11v+10=0$, so $(v-1)(v-10)=0$. The root $v=1$ would make the second speed negative, so $v=\boxed{10}$ km per hour.

Check: $\dfrac{5}{10}+\dfrac{4}{8}=\dfrac{1}{2}+\dfrac{1}{2}=1$ hour.
